\section{Two raw positive realizations}\label{sec:realizations}
\subsection{Classical relative transport}
Fix $d\ge1$, $k=d+1$, a known point $x_*\in\mathbb S^d$, and a known contrast $0<\eta\le1$. One raw realization has hidden physical state a finite word $w$ of rotations, in the standard Borel disjoint union $\bigsqcup_{j\ge0}SO(k)^j$. Preparation creates the empty word. A call appends an independent Haar rotation and reports that rotation, but never allows rereading the word. Set $x(w)$ to its ordered product applied to $x_*$. The induced coordinate kernel is
\begin{equation}\label{eq:classical-kernel}
 K(dx',dR\mid x)=\delta_{Rx}(dx')\,h(dR),
\end{equation}
where $h$ is Haar probability on $SO(k)$. Only $R$ is reported. The audit chooses $I$ uniformly from $\{1,\ldots,k\}$ and then $S\in\{-1,1\}$ with probability $(1+S\eta x_I)/2$. Its scored vector is $Z=(k/\eta)S e_I$, so $\E[Z\mid x]=x$ and $\norm Z^2=k^2/\eta^2$. The audit happens after the answer has been committed and is not a further observation supplied to the machine.

\begin{theorem}[Classical realization]\label{thm:classical}
For the raw preparation, kernel~\eqref{eq:classical-kernel} and single-audit interface, the predictive quotient is the unit sphere with its Borel structure, and the hypotheses of Theorem~\ref{thm:energy} hold with $A(R)=R$, $\kappa=1$ and $B=k^2/\eta^2-1$. Theorems~\ref{thm:sphere} and~\ref{thm:clock-circle} therefore give the sequential and autonomous resource curves; on the circle Theorem~\ref{thm:circle-product} is exact. Suppressing reports gives exactly the spherical--atomic law and estimates of Theorem~\ref{thm:singular}.
\end{theorem}
\begin{proof}
The word kernel is $w\mapsto w\mathbin{\|}R$ with law $h(dR)$. Thus no finite-dimensional manifold was assumed for the raw state. Nonnegativity and normalization of its induced coordinate kernel follow from the two factors and $|\eta x_I|\le1$. Conditional on a reported $R$, the coordinate functions pull back to $Rx$. Their squared integral matrix is $I$. The report law is independent of every entering point and all earlier histories, so the required conditional independence is an identity of the raw kernel, not a filtering approximation. Coordinate audits separate all points. Any finite continuation of reported rotations followed by the one audit remains affine in the entering coordinates, including report-dependent choices of the final audit, because the report law is state independent. Thus the finite test space closes. The Borel section described in Section~\ref{sec:experiments} proves realized minimality.

A performed first Haar acquisition sends $x_*$ to $\sigma_d$, and each subsequent acquisition preserves that law; this proves the actual acquired measure. The update is an isometry on the entire sphere and unit ball, and the report law has total-variation distance zero between different predictive points. The audit probabilities are Lipschitz in the coordinates. Thus global covering, small-ball mass and update stability are precisely the conditions proved in Lemma~\ref{lem:sphere-geometry}, not supplied model-specific exponents. Conditional on a hidden Haar transform the coordinate mean is zero, and subsequent visible rotations preserve it, proving the singular claim. All response tables and rotations used by the observer are independent of the actual hidden point. There are exactly $N=n+2$ physical calls.
\end{proof}
The full predictive quotient here refers to the declared future experiment, with one terminal linear audit. Adding repeated nonlinear observations changes the test space; it is not covered merely by retaining the same coordinates. Known orthogonal actions selected from the counted label do not improve the lower: conditioning on their fresh randomness leaves the reported Haar product isotropic. Actions that reveal the absolute point, undo an unknown transform from a free history tape, or reprepare the system are not in this interface.

\subsection{A noncommuting quantum instrument}
Let $\sigma_1,\sigma_2,\sigma_3$ be the Pauli matrices. Prepare
\[
 \rho_*=(I+r x_*\cdot\sigma)/2,\qquad x_*\in\mathbb S^2,\quad0<r\le1.
\]
The contrast $r$ is a specified preparation parameter, not a value learned for free. For Haar probability $h$ on $SU(2)$, the reported positive instrument is
\begin{equation}\label{eq:quantum-kernel}
 \mathcal I(dU)(\rho)=U\rho U^*\,h(dU).
\end{equation}
Each density map is completely positive, and the integrated instrument is trace preserving. Its report probability is $h(dU)\tr\rho$. The audit chooses $I\in\{1,2,3\}$ uniformly and measures the two-outcome effect $(I+S\sigma_I)/2$. The scored vector is $Z=(3/r)S e_I$. No joint measurement of the three Pauli observables is asserted.

\begin{theorem}[Noncommuting realization and rank change]\label{thm:quantum}
For the instrument~\eqref{eq:quantum-kernel}, the test coordinate $x$ in $\rho=(I+r x\cdot\sigma)/2$ is a realized predictive quotient for the declared instrument--audit interface. The same general energy and sphere theorems apply with $d=2$, $\kappa=1$, $V=9/r^2$. The state family and the audit effects do not share an eigenbasis. Replacing the instrument by
\begin{equation}\label{eq:quantum-erased}
 \mathcal I_t(dU)(\rho)=\alpha_t U\rho U^*h(dU),\qquad
 \mathcal I_t(\{\bot\})(\rho)=(1-\alpha_t)\tr(\rho)I/2
\end{equation}
verifies Theorem~\ref{thm:singular}. For $r=1$, the actual conditional states change from rank one to rank two on the first hidden report. The true rank-one mass is $w_t$, not one.
\end{theorem}
\begin{proof}
The identities $U(x\cdot\sigma)U^*=(R_Ux)\cdot\sigma$ define the usual surjection from $SU(2)$ to $SO(3)$. The pushforward of Haar probability is Haar probability because it is normalized and invariant. The instrument therefore transports the centered test coordinates by $A(U)=R_U$ with state-independent reports and $\int A^TA\,dh=I$. The three single-query effects separate all Bloch coordinates; every allowed future continuation remains affine in $\rho$. A Borel choice of a unitary taking $x_*$ to $x$ on two coordinate charts gives a section. The audit calculation yields $\E[Z\mid\rho]=x$ and $\E\norm Z^2=9/r^2$.

For two nonparallel vectors $x,y$, the commutator is proportional to $i r^2(x\times y)\cdot\sigma$, and is nonzero; thus a shared diagonal representation would be false. All continuity, mass and stability requirements nevertheless follow from the same spherical calculations, since the report probability does not depend on $x$. Haar twirling has $\int U\rho U^*dh=\tr(\rho)I/2$, obtained by invariance and trace. The hidden branch of~\eqref{eq:quantum-erased} is therefore the actual conditional state after an unreported unitary, not an artificial reset inserted by the observer. Once it occurs, later unitaries preserve $I/2$. This proves the mixed law and the rank statement, and hence all the singular inequalities.
\end{proof}
The physical qubit is the experiment being predicted, not uncharged observer memory: the observer may neither store an ancillary quantum state nor perform extra measurements. Only its classical labels survive between reports. Completely positive finite-dimensional matrix identities are all that is used; no conclusion about unbounded-operator closures follows. Both realizations verify the hypotheses of the general theorem after that theorem has been proved, not conversely.

\subsection{Higher-dimensional pure-state orbits: intrinsic rather than ambient dimension}
Let $D\ge2$, $V_D$ be the real Hilbert space of traceless Hermitian $D\times D$ matrices with the Hilbert--Schmidt norm, and $c_D=\sqrt{1-1/D}$. Prepare a known rank-one projector $\Pi_*$ and use Haar conjugation on $SU(D)$. The centered coordinate is
\[
 X=(\Pi-I/D)/c_D,\qquad
 \mathcal O_D=\{(\Pi-I/D)/c_D:\Pi\text{ rank one}\}.
\]
Its norm is one, but its acquired dimension will be $d_D=2(D-1)$, not $\dim V_D=D^2-1$. These facts will follow from the raw unitary law.

\begin{lemma}[Uniform mass across all conjugacy orbits]\label{lem:quantum-orbits}
For every $D\ge2$ there are constants $C_D,r_D>0$ such that for every traceless Hermitian $V$ with $\norm V_{\rm HS}=1$, every Hermitian center $Z$ and $0<r\le r_D$,
\begin{equation}\label{eq:unitary-smallball}
 h\{U:\norm{UVU^*-Z}_{\rm HS}\le r\}\le C_D r^{2(D-1)}.
\end{equation}
The pure-state orbit $\mathcal O_D$ has global covers of radius $C_Dm^{-1/(2(D-1))}$. The stabilizer of any one of its points has fixed subspace exactly its line in $V_D$.
\end{lemma}
\begin{proof}
Write the eigenvalues of $V$ in decreasing order. Their sum is zero and sum of squares is one, so their range is at least $1/\sqrt D$. Some consecutive gap is at least $g_D=1/(\sqrt D(D-1))$, say between ranks $j$ and $j+1$. For two matrices $W,W'$ in this orbit, let $P,P'$ be their top-$j$ spectral projections. The variational characterization of the top eigenspace gives
\[
 \tr(W(P-P'))\ge g_D(j-\tr(PP')),
 \quad \tr(W'(P'-P))\ge g_D(j-\tr(PP')).
\]
Adding and using Cauchy--Schwarz yields
$\norm{P-P'}_{\rm HS}\le g_D^{-1}\norm{W-W'}_{\rm HS}$.
If the ball in~\eqref{eq:unitary-smallball} meets the orbit, choose one $W'$ in it; all other points in the ball have their top projections within $2r/g_D$ of $P'$. The random top projection is Haar on the complex Grassmannian of rank $j$.

For clarity its local mass estimate follows directly from Haar frames. Near a fixed rank-$j$ plane use the graph chart $A\in\mathbb C^{(D-j)\times j}$. Normalizing a full-rank complex Gaussian frame gives Haar measure. Write that frame as $(B^T,C^T)^T$ with square invertible upper block $B$, and substitute $C=AB$. The real Jacobian is $|\det B|^{2(D-j)}$. Integrating $\exp[-\tr B^*(I+A^*A)B]|\det B|^{2(D-j)}$ over $B$, then substituting $(I+A^*A)^{1/2}B$, gives density proportional to $\det(I+A^*A)^{-D}$. On a fixed neighborhood this density is bounded and projector distance is comparable to $\norm A_{\rm HS}$. Thus a radius-$s$ ball has mass at most $C_{D,j}s^{2j(D-j)}$. Unitary invariance makes the constant independent of the center. Since $2j(D-j)\ge2(D-1)$ and there are only $D-1$ possible ranks, the stated uniform bound follows. This proof does not assume a uniform simple spectrum.

For rank one, the same chart density is also bounded below locally and the chart has $2(D-1)$ real coordinates. Compactness, finitely many charts and a maximal separated set give the asserted global covers. Finally the stabilizer of a rank-one projector acts arbitrarily unitarily on its orthogonal complement (with the determinant corrected on the first line). Its fixed Hermitian matrices are scalar on each of the two blocks; the traceless condition leaves exactly the line through $\Pi-I/D$.
\end{proof}

\begin{theorem}[Projective quantum acquired-geometry law]\label{thm:projective}
Prepare $(1-r)I/D+r\Pi_*$ with specified $0<r\le1$ and use the raw instrument $\mathcal I(dU)(\rho)=U\rho U^*h(dU)$. Let $F_1,\ldots,F_k$ be any Hilbert--Schmidt orthonormal basis of $V_D$, $k=D^2-1$. At audit draw $I$ uniformly and measure the effects $(I+S F_I)/2$, which are positive since $\norm{F_I}_{\rm op}\le1$. Score $Z=(k/(r c_D))S e_I$. Then the hypotheses of Theorem~\ref{thm:homogeneous} hold with
\[
 d=2(D-1),\quad V=k^2/(r^2c_D^2),\quad B=V-1.
\]
In particular, for $M\ge n+1$,
\begin{equation}\label{eq:projective-rate}
 \calR_{\rm aut}(n,M)-B\asymp_D
 \min\{1,n^{D/(D-1)}(M-1)^{-1/(D-1)}\},
\end{equation}
whereas the checkpoint excess is of order $m^{-1/(D-1)}$. Independent suppressed reports verify Corollary~\ref{cor:homogeneous-singular}; for $r=1$ the conditional rank changes from one to $D$ with actual masses $w_n,1-w_n$.
\end{theorem}
\begin{proof}
The centered raw audit mean is exactly the coordinate vector of $X=(\rho-I/D)/(r c_D)$, and its second moment is the displayed $V$. Conjugation acts orthogonally on $V_D$; its only invariant traceless matrix is zero. The word of reported unitary instruments pulls back the finite coordinate tests linearly. As in the qubit proof, these tests separate the declared conditional states and have a Borel section, with pure-state projectors constructed from local normalized vectors. Haar conjugation derives the actual orbit measure. Lemma~\ref{lem:quantum-orbits} verifies both the mass certificate for every possible incoming mean direction and the global cover and stabilizer certificates. This uniformity is crucial: incoming means need not themselves lie on the pure-state orbit. Apply the general theorem and substitute $d=2(D-1)$ to obtain~\eqref{eq:projective-rate}. Haar twirling sends every state to $I/D$, proving the hidden-report and rank assertions. The instrument family is noncommuting for nonorthogonal pure-state preparations, and no shared-basis reduction was used.
\end{proof}
